% ============================================================================
% 第四章：总结与展望
% ============================================================================

\section{研究总结}

\begin{frame}{研究背景与问题}
    \begin{itemize}
        \item \textbf{MI-BCI 性能瓶颈}: 时间窗选择是关键因素
        \item \textbf{现有方法局限}:
        \begin{itemize}
            \item 固定时间窗：无法适应个体间差异
            \item 贝叶斯优化：计算复杂度高
            \item 深度强化学习：模型复杂、训练不稳定、可解释性差
        \end{itemize}
        \item \textbf{核心矛盾}: 深度强化学习对大量训练数据和高计算资源的依赖 \\
        vs \\
        BCI 系统中数据稀缺、设备资源受限的现实
    \end{itemize}
\end{frame}

\begin{frame}{研究方法}
    \begin{block}{基于表格型 Q-Learning 的轻量化时间窗优化方法}
        \begin{itemize}
            \item 将时间窗优化建模为 MDP (136 个离散状态)
            \item 使用经典时序差分学习迭代优化策略
            \item 以 CSP-SVM 分类准确率为奖励信号
        \end{itemize}
    \end{block}
    
    \vspace{0.3cm}
    \begin{exampleblock}{核心优势}
        \begin{itemize}
            \item \textbf{轻量化}: 544 个参数 (vs DQN 百万级)
            \item \textbf{稳定性}: 理论收敛保证
            \item \textbf{可解释性}: Q 值表可直接可视化
        \end{itemize}
    \end{exampleblock}
\end{frame}

\begin{frame}{主要实验结果}
    \begin{table}[h]
        \centering
        \begin{tabular}{lcc}
            \hline
            \textbf{指标} & \textbf{Table Q} & \textbf{DQN} \\
            \hline
            准确率 & 62.68\% & 62.70\% \\
            统计检验 & \multicolumn{2}{c}{$p = 0.981$} \\
            训练时间 & \textbf{快 3 倍} & 慢 \\
            推理延迟 & \textbf{低 1000 倍} & 高 \\
            内存占用 & \textbf{低 125 倍} & 高 \\
            \hline
        \end{tabular}
    \end{table}
    
    \vspace{0.3cm}
    \begin{alertblock}{核心结论}
        表格型 Q-Learning 在保持与 DQN 相当性能的同时，\\
        显著降低了计算资源需求
    \end{alertblock}
\end{frame}

% ============================================================================
% 主要贡献
% ============================================================================
\section{主要贡献}

\begin{frame}{理论贡献}
    \begin{enumerate}
        \item \textbf{轻量化 BCI 设计新范式}
        \begin{itemize}
            \item 将深度强化学习简化为表格型强化学习
            \item 实现轻量化与高性能的统一
        \end{itemize}
        
        \item \textbf{算法选择准则}
        \begin{itemize}
            \item 状态空间 $< 10^4$ 且可充分访问时
            \item 表格型方法在收敛性、样本效率、可解释性上优于 DQN
        \end{itemize}
        
        \item \textbf{可解释性 BCI 系统实践}
        \begin{itemize}
            \item Q 值表可直接可视化和分析
            \item 为医疗康复等安全敏感应用提供支撑
        \end{itemize}
    \end{enumerate}
\end{frame}

% ============================================================================
% 局限性
% ============================================================================
\section{局限性}

\begin{frame}{研究局限性}
    \begin{itemize}
        \item \textbf{验证范围局限}
        \begin{itemize}
            \item 仅在 9 名健康被试的标准数据集上验证
            \item 未在真实患者群体或在线 BCI 系统中测试
        \end{itemize}
        
        \item \textbf{方法设计局限}
        \begin{itemize}
            \item 时间窗离散化步长 0.2 秒，可能无法定位全局最优
            \item 奖励信号为间接辅助任务 (CSP-SVM 准确率)
        \end{itemize}
        
        \item \textbf{通用性局限}
        \begin{itemize}
            \item 仅针对 CSP-SVM 分类器设计
            \item 未在其他 BCI 范式 (SSVEP, P300) 上验证
        \end{itemize}
    \end{itemize}
\end{frame}

% ============================================================================
% 未来展望
% ============================================================================
\section{未来展望}

\begin{frame}{未来工作方向}
    \begin{enumerate}
        \item \textbf{在线和临床验证}
        \begin{itemize}
            \item 移植到嵌入式 BCI 系统
            \item 在脑卒中康复患者群体中验证
        \end{itemize}
        
        \item \textbf{优化算法效率}
        \begin{itemize}
            \item 基于元学习的快速奖励估计
            \item 缓存机制复用交叉验证结果
        \end{itemize}
        
        \item \textbf{探索泛化能力}
        \begin{itemize}
            \item 与深度学习模型 (EEGNet) 结合
            \item 扩展到通道选择、CSP 滤波器优化
            \item 时空联合优化
        \end{itemize}
        
        \item \textbf{迁移学习策略}
        \begin{itemize}
            \item 提取通用时间窗选择先验知识
            \item 新用户仅需 10--20 个试次快速校准
        \end{itemize}
    \end{enumerate}
\end{frame}

% ============================================================================
% 致谢
% ============================================================================
\section{致谢}

\begin{frame}{致谢}
    \begin{center}
        \Large
        感谢指导老师的悉心指导！\\
        \vspace{0.5cm}
        感谢评阅老师在百忙之中评阅本论文！\\
        \vspace{0.5cm}
        感谢各位同学的聆听！
    \end{center}
\end{frame}

\begin{frame}
    \begin{center}
        {\Huge Q\&A}
    \end{center}
\end{frame}
